\begingroup
\setlength{\parskip}{.25em}
\section*{\personaCuatro · Llamadas y fallos de red}
\textbf{Proyecto}\quad \proyectoCuatro

\subsection*{1. Conceptos y explicación}
\texttt{send()} mantiene esperando al hilo que llama. \texttt{sendAsync()} devuelve un futuro y permite continuar. En la demo, \texttt{handle()} procesa el resultado y \texttt{join()} espera al final para conservar la consola. Ambos modos necesitan gestionar el timeout y los errores.

\subsection*{2. Demostración Java}
Byron ejecutó los dos programas de \texttt{demo/opcion4/} el 20 de noviembre. Las primeras consultas normales devolvieron HTTP 200 en 149 y 77 ms. El retardo controlado produjo \texttt{HttpTimeoutException} en 3021 y 3027 ms. Con el servidor detenido, ambos clientes reportaron conexión fallida en 52 y 61 ms. Tras reiniciarlo, otra consulta devolvió 200 en 159 ms.

El futuro estaba pendiente cuando continuó el hilo. Los tiempos son observaciones individuales y no demuestran que un modo sea más rápido. El código limita la conexión a 2 s y la petición a 3 s. \texttt{/lenta} demora 5 s. Un timeout significa que el cliente dejó de esperar, sin probar que el servidor cancelara su trabajo.

\subsection*{3. Aplicación al PFC CLAB}
El esqueleto tiene dos procesos, Reservas y Laboratorios. \texttt{GET} \path{/api/reservas/comunicacion} consulta \texttt{GET} \path{/api/laboratorios/info} y devuelve la información recibida. La llamada interna con \texttt{RestClient} es bloqueante. El cliente de demostración permite elegir ambos modos. El código heredado de gestión de reservas está en \texttt{legado/}.

\subsection*{4. Decisiones distribuidas}
La API oculta las clases Java mediante JSON. Las direcciones y los fallos quedan visibles. La prueba demuestra un fallo de proceso, no una partición de red.

Para una futura confirmación de reservas replicada, se propone impedir que una instancia aislada confirme un horario sin acuerdo con la autoridad de escritura. Durante una partición, esto sacrifica disponibilidad de confirmación para conservar consistencia. Sin partición, la propuesta prioriza validar y registrar el horario con datos consistentes, aceptando la coordinación y su latencia. Son decisiones de diseño guiadas por CAP/PACELC, no garantías implementadas ni medidas por este esqueleto.

\subsection*{5. Fallo, recuperación y conclusión}
Byron detuvo Laboratorios. La consulta interna produjo una excepción del cliente HTTP y Reservas respondió 503, recibido en 178 ms. Sus límites de conexión y lectura son de 2 s cada uno. Reservas siguió atendiendo peticiones, pero su consulta dependiente quedó indisponible. No hubo failover ni reintentos automáticos. Tras el reinicio manual de Laboratorios, otra consulta obtuvo 200 en 312 ms. Un intento con la ruta mal copiada terminó en \texttt{comunicacionjava} y recibió 404.

MRBS ofrece reservas de espacios y prevención de conflictos. El esqueleto de CLAB aún no implementa esa operación de negocio. Su separación permitió reiniciar Laboratorios sin reiniciar Reservas, con el costo de gestionar una dependencia de red. La elección se justifica si esa independencia compensa la complejidad operativa. Para futuras escrituras, un timeout puede dejar incierto el resultado. Antes de reintentar se propone identificar la operación para evitar duplicados, limitar los intentos y espaciar las nuevas peticiones.

\subsection*{6. Caso documentado de apoyo: GitHub}
El 14 de agosto de 2024, una configuración errónea a las 22:59 UTC impidió responder a comprobaciones de salud. Los hosts quedaron fuera del enrutamiento y las lecturas críticas resultaron inaccesibles. GitHub.com estuvo indisponible de 23:02 a 23:38 UTC. Revirtieron el cambio, verificaron la conectividad y continuaron observando hasta las 00:30 UTC del día 15. No reportaron pérdida ni corrupción de datos. Añadieron controles sobre cambios y priorizaron reversión y resiliencia.

La cadena separa el detonante de su propagación. El informe no identifica el parámetro exacto cambiado. La lección propuesta para CLAB es comprobar dependencias y preparar reversión. Este caso pertenece a GitHub y no es un incidente de CLAB.

\subsection*{7. Evidencias y fuentes}
Las transcripciones del ensayo están en \path{integrante-4/evidencias/demo-http/}. Se conservaron los comandos, los mensajes y los tiempos aportados por Byron.
\par
{\small Fuentes de apoyo\quad
\href{https://docs.oracle.com/en/java/javase/17/docs/api/java.net.http/java/net/http/HttpClient.html}{Oracle Java 17},
\href{https://github.com/Danxxx777/clab-pfc-7a/tree/68cf3dbf1ddd5133ca8e1b430cfb113cfd10341d/src}{código de CLAB},
\href{https://doi.org/10.32620/reks.2023.2.14}{Ahmed et al. (2023)},
\href{https://meeting-room-booking-system.github.io/}{MRBS},
\href{https://learn.microsoft.com/en-us/azure/architecture/patterns/retry}{Microsoft} y
\href{https://github.blog/news-insights/company-news/github-availability-report-august-2024/}{informe de GitHub (2024)}.
}
\IfFileExists{../evidencias/esquema-mano/clab-bosquejo.jpg}{%
\begin{center}
\includegraphics[width=.85\linewidth,height=35mm,keepaspectratio]{../evidencias/esquema-mano/clab-bosquejo.jpg}
\end{center}
}{}
\endgroup
